% Software requirements (input file of touchsynth_poly_report.tex). Draft for review, 2026-10-20.
\newcommand{\stmet}{\textcolor{green!40!black}{\textbf{met}}}
\newcommand{\stpart}{\textcolor{orange!80!black}{\textbf{partial}}}
\newcommand{\stopen}{\textcolor{red!70!black}{\textbf{open}}}
\newcommand{\stdes}{\textcolor{blue!60!black}{\textbf{design}}}
\newcommand{\review}[1]{\par\smallskip\noindent\fcolorbox{blue!60!black}{blue!5}{\parbox{0.97\linewidth}{%
  \small\textbf{For review.} #1}}\par\smallskip}

% =====================================================================================================================
\section{Software requirements}\label{sec:srs}

This section states what the software of the instrument must do, as requirements that can be verified one by one.
They are derived from the system requirements (Tables~\ref{tab:reqf} and~\ref{tab:reqp}; \file{docs/design.md} 2), the voice bus protocol
(\file{docs/bus.md}), the MIDI and patch model (\file{docs/midi.md}) and the controller plan
(\file{docs/controller.md}). Section~\ref{sec:sad} describes the architecture that implements them, and
Section~\ref{sec:fw} describes the firmware as built.

\review{This is a first draft, status 20 October 2026. Please check three things: the scope of each requirement
(Section~\ref{sec:srsdec} lists the choices that are still open); the numeric targets marked \emph{proposed}; and the
split between what the cards do and what the controller does.}

\subsection{Scope, terms and conventions}

\paragraph{Software items.}
\begin{itemize}
  \item \textbf{Card firmware (CF):} \file{firmware/voice}, built with \texttt{make BOARD=card} for the CH32V006 on each
    voice card. The same code base builds for board~1, the rev 0.1 test bed, which keeps its analog DCO and its MIDI
    input.
  \item \textbf{Controller firmware (CT):} the backplane's CH32V203. Today \file{firmware/bench\_ctrl} stands in for it
    on a development board. It implements the frame engine, the bus, USB-MIDI and CDC, and the SWIO bridge, but not
    the parts, the storage or the 74HC595.
  \item \textbf{Voice bus (BUS):} the protocol between the two firmwares. Its requirements bind both sides.
  \item \textbf{Tools (TL):} the host-side generators, tests and bench scripts (\file{tools/},
    \file{firmware/voice/test/}, \file{firmware/voice/gen\_tables.py}, \file{sim/behav/}).
\end{itemize}

\paragraph{Columns.} \emph{Source} names the system requirement or document the requirement comes from.
\emph{Verification} names the method:
\begin{itemize}
  \item T: an automated test, named;
  \item B: a bench measurement;
  \item I: an inspection or review;
  \item S: a simulation or model;
  \item H: needs the new boards (verification stage V2--V4, Table~\ref{tab:vstages}).
\end{itemize}
\emph{Status} says how far the requirement is shown today:
\begin{itemize}
  \item \stmet: shown on the bench (CH32V203 development board plus one V006 card) or on board~1;
  \item \stpart: implemented in part, or shown only in a bench variant;
  \item \stopen: not implemented;
  \item \stdes: follows from the design, to be confirmed on hardware.
\end{itemize}

\paragraph{Terms.} A frame is the \SI{1}{ms} between two rising SYNC edges. $L$ is the bus lead, two frames; $D$
is the presentation delay, four frames. A buffer half is 256 carrier periods of the PWM DAC (\SI{1.37}{ms}). A carrier
period is \SI{5.33}{\micro s}. A patch slot is one of the 4 patches a card holds. A part is one instrument of the
multitimbral setup.

% ---------------------------------------------------------------------------------------------------------------------
\subsection{Card firmware}\label{sec:srscf}

\begingroup\small\setlength{\tabcolsep}{4pt}
\begin{longtable}{@{}>{\raggedright}p{0.9cm} >{\raggedright}p{7.5cm} >{\raggedright}p{1.5cm}
  >{\raggedright}p{2.7cm} >{\raggedright\arraybackslash}p{2.6cm}@{}}
\caption{Card firmware requirements.}\label{tab:srscf}\\
\toprule
ID & Requirement & Source & Verification & Status \\
\midrule\endfirsthead
\toprule
ID & Requirement & Source & Verification & Status \\
\midrule\endhead
\bottomrule\endfoot
\multicolumn{5}{@{}l}{\textit{Start-up, time base and clock}} \\
CF1 & From reset until the first SYNC edge, and indefinitely without SYNC, the card runs standalone: SysTick
  ticks, a mono synth on its built-in presets, the SWIO console, the receiver off. The first SYNC edge starts the
  trim (CF3). & F8, bus.md 6 & T \texttt{sync}; B & \stmet \\
CF2 & The control tick runs on each rising SYNC edge. Edges less than $3/4$ frame after the last are ignored and
  counted. A missing edge is bridged by a SysTick tick \SI{0.25}{ms} late. After 10 frames without SYNC the VCA is
  muted, and the card re-locks when SYNC returns. & P10, bus.md 6 & T \texttt{sync} (gaps of 3 and 40 frames) &
  \stmet \\
CF3 & After 100 SYNC periods the card sets HSITRIM so that its HSI is within \SI{\pm0.1}{\%} of the crystal. It
  enables its receiver only after that trim. & P1, bus.md 2, 6 & B (STATUS clock error 186--617\,ppm) & \stmet \\
CF4 & The oscillator phase steps are scaled by the measured SYNC period, filtered over about 64 frames, so the
  pitch follows the controller's crystal and voices cannot drift apart. & P1, P11 & B; H (V3: two cards against the
  crystal) & \stpart{} (built; error against the crystal not measured) \\
\multicolumn{5}{@{}l}{\textit{Voice bus, card side}} \\
CF5 & The card receives packets at \SI{1}{Mbaud} by DMA into a ring and parses them in the control tick: COBS,
  CRC-8. A packet for another card is dropped at its dst byte, before the CRC. Every error class is counted. &
  bus.md 4, 9.5 & T \texttt{crc}, \texttt{load} & \stmet \\
CF6 & A timed event runs in the tick of its execution frame. A late event runs at once and is counted. An event
  more than 32 frames ahead is dropped and counted. At least 16 events can wait. & bus.md 4 & T \texttt{late} &
  \stmet \\
CF7 & Whatever a tick computes reaches the outputs $D$ frames later: the CVs through a delay line, the DAC state
  exact to one carrier period. The VCA trim is not delayed. & P8, P10, bus.md 3.1 & B (scope: 4.008\,ms, spread
  \SI{8}{\micro s}) & \stmet \\
CF8 & An unaddressed card answers ASSIGN to dst 0x00 with HELLO: its 64-bit chip ID, firmware version and
  capabilities. After a reset over NRST the card is found again within \SI{0.4}{s}. & F1, bus.md 6 & T
  \texttt{reset}; H (V3: NRST from the 74HC595) & \stpart{} (shown with an SWIO reset; the NRST path needs the
  backplane) \\
CF9 & The card answers REQUEST (STATUS, COUNTERS, HELLO) within 2 frames, on an open-drain return line that it
  releases when idle. & bus.md 2, 5.2 & T \texttt{enum}, \texttt{counters}; B (open drain with \SI{1}{k\ohm}) & \stpart{} (the bench
  ran the push-pull variant, \texttt{BUS\_TX\_PP}; the open-drain build is untested) \\
\multicolumn{5}{@{}l}{\textit{Voice engine and sound}} \\
CF10 & One voice per card at a \SI{1}{kHz} control rate: amplitude and filter envelopes, LFO, glide, pitch envelope,
  and the cutoff, resonance and VCA laws onto the CVs. & F5 & B (board~1, bench) & \stmet \\
CF11 & DSP engines: up to 3 band-limited oscillators (saw and pulse, polyBLEP), noise, and modal drums in fixed point,
  one sample per \SI{187.5}{kHz} carrier period. The output is a 7-bit PWM DAC with second-order error feedback. The
  render is bit-exact with the previous revision and follows \file{sim/behav/dsp.py}. & F5, P2 & T
  \texttt{render\_test.py} (197 scenarios), card \texttt{check}; S (S9) & \stmet \\
CF12 & No DAC buffer half is missed with any patch under 8-voice bus load or the stress test. Each half keeps at least
  50 carrier periods of slack (\emph{proposed} margin). & F5 & T \texttt{load --voices 8}, \texttt{stress} &
  \stmet{} (least 99 periods) \\
CF13 & Four CVs as 10-bit PWM with 2-bit dither at \SI{46.875}{kHz}, loaded by DMA without a CPU interrupt, on
  TIM1 CH1--CH4 (PC6, PC7, PC0, PD3). & F5 & B (pin activity on the V006) & \stmet \\
CF14 & Stereo pan: a 7-bit PWM on PC5 (TIM2 CH1, the DAC's carrier) sets the steering pairs. Pan law: linear
  steering, with CV\_VCA raised towards the centre for constant power and capped at 70.6\,\% at hard pan. Pan is set
  per voice by PARAM. & F7, spec 3.1 & B (PC5 output shown with TIM2 CH1 enabled by hand); H (V2: L / R levels) &
  \stopen{} (pin and timer shown; firmware missing) \\
CF15 & Patches: 4 patch slots per card. PATCH\_DATA writes a slot that is not active; a timed PATCH\_SELECT switches
  on the tick. All parameters of midi.md 4 are implemented. & F6, midi.md 4 & T \texttt{patch}; I & \stpart{}
  (built-in presets through PATCH\_PRESET; PARAM IDs 20, 28, 32--35, 38--47 only) \\
CF16 & Performance controls (CTRL): bend with the patch's BEND\_RANGE, mod wheel, channel pressure, expression,
  tempo, LFO phase reset. & midi.md 3, bus.md 5.1 & T \texttt{ctrl} & \stpart{} (bend fixed at \SI{\pm2}{}
  semitones; pressure and tempo missing) \\
\multicolumn{5}{@{}l}{\textit{Calibration and self-test}} \\
CF17 & On CALIBRATE the card mutes, measures cutoff, the resonance threshold and the VCA trim through its loopback ADC
  (PC4) and, if fitted, the dual-PWM weight. It reports busy in STATUS and delivers its results as CAL\_DATA. & F9,
  bus.md 6 & H (V2) & \stopen{} (rev 0.1 calibration on board~1; the \SI{5}{V} tables need a card) \\
CF18 & Self-test inputs (option): the ladder bias on PA2 (A0) and the analog rail on PD2 (A3), readable over the
  console and STATUS. & spec 3.2 & H & \stopen{} (option) \\
\multicolumn{5}{@{}l}{\textit{Maintenance, resources and safety}} \\
CF19 & The card is programmed and debugged over SWIO alone, by a WCH-LinkE or the controller (minichlink compatible),
  and has no bootloader of its own. PD7 stays configured as NRST (option byte RST\_MODE $\neq$ 11). & F4, F8 & B
  (40\,s through the V203 bridge; option bytes 0xF7 read) & \stmet \\
CF20 & SWIO console: status, counters, the DAC and render diagnostics (\texttt{dac}, \texttt{bench}, \texttt{check}),
  ADC captures, manual CVs. & F8 & B & \stmet \\
CF21 & Resources: static RAM at most \SI{7.2}{KB} of \SI{8}{KB}; at least 512 bytes of stack never used under bus
  load; flash at most \SI{60}{KB} of \SI{62}{KB} (\emph{proposed} limits). & C3 & B (stack low-water mark), I (link
  map) & \stmet{} (\SI{7.0}{KB}, 656 bytes, \SI{55}{KB}) \\
CF22 & PA4 is never configured as an output: it is bonded to PD7 (NRST) in the TSSOP-20. & C3 & I (code review) &
  \stmet \\
CF23 & Build options: dual-PWM DAC (PC1, TIM2 CH4) and an external \SI{24}{MHz} clock on PA1. & spec 3.5, 6 & B
  (board~1); H & \stpart{} (the dual-PWM build compiles; untested on a V006 card) \\
\end{longtable}
\endgroup

% ---------------------------------------------------------------------------------------------------------------------
\subsection{Voice bus}\label{sec:srsbus}

\begingroup\small\setlength{\tabcolsep}{4pt}
\begin{longtable}{@{}>{\raggedright}p{0.9cm} >{\raggedright}p{7.5cm} >{\raggedright}p{1.5cm}
  >{\raggedright}p{2.7cm} >{\raggedright\arraybackslash}p{2.6cm}@{}}
\caption{Voice bus requirements (both firmwares).}\label{tab:srsbus}\\
\toprule
ID & Requirement & Source & Verification & Status \\
\midrule\endfirsthead
\toprule
ID & Requirement & Source & Verification & Status \\
\midrule\endhead
\bottomrule\endfoot
BUS1 & \SI{1}{Mbaud} 8N1 on BCAST and RETURN; COBS framing with 0x00 delimiters; CRC-8/SMBUS; at most 64 decoded
  bytes per packet; little-endian values; unknown message types ignored and counted. & bus.md 2, 4 & T \texttt{all}
  & \stmet \\
BUS2 & Every frame starts with FRAME at the SYNC edge. A card's frame counter that disagrees with FRAME is set again
  and counted as a slip. & bus.md 3 & T \texttt{sync} & \stmet \\
BUS3 & Timed events execute at the frame $a + L$ of their arrival; sound follows $D$ frames later, the same for
  every event and every voice. USB-MIDI to sound at most \SI{10}{ms}, constant. & P8 & B (\texttt{latency
  --scope}: 5.0--6.0\,ms) & \stmet \\
BUS4 & The same event sounds on every card within \SI{10}{\micro s}. & P10 & B (one card: spread \SI{8}{\micro s});
  H (V3/V4: two or more cards) & \stdes \\
BUS5 & Broadcast without acknowledgement. A VOICE\_STATE refresh repairs a lost NOTE\_ON or NOTE\_OFF within
  \SI{50}{ms}; a resent patch repairs a lost PARAM. & bus.md 1, 5.3 & T \texttt{repair} & \stpart{} (patch resend
  by PATCH\_PRESET only) \\
BUS6 & One card speaks at a time, only when asked: one REQUEST or ASSIGN outstanding; the controller times out after
  5 frames. & bus.md 5.2 & T \texttt{enum}, \texttt{reset} & \stmet \\
BUS7 & An 8-voice chord fits one frame. Per-voice records and notes ahead of controllers keep legato chord changes and
  multi-part controllers inside the bus lead. & bus.md 5.3, 9.5 & T \texttt{load} & \stmet \\
BUS8 & Enumeration of 8 slots and power-up to playable within \SI{3}{s} (P9): about \SI{1}{s} plus \SI{0.3}{s} per
  empty slot. & F1, P9 & H (V3/V4) & \stdes{} (one card: 0.3--0.5\,s) \\
\end{longtable}
\endgroup

% ---------------------------------------------------------------------------------------------------------------------
\subsection{Controller firmware}\label{sec:srsct}

Status refers to \file{firmware/bench\_ctrl} where it implements the function; the backplane build does not exist yet.

\begingroup\small\setlength{\tabcolsep}{4pt}
\begin{longtable}{@{}>{\raggedright}p{0.9cm} >{\raggedright}p{7.5cm} >{\raggedright}p{1.5cm}
  >{\raggedright}p{2.7cm} >{\raggedright\arraybackslash}p{2.6cm}@{}}
\caption{Controller firmware requirements.}\label{tab:srsct}\\
\toprule
ID & Requirement & Source & Verification & Status \\
\midrule\endfirsthead
\toprule
ID & Requirement & Source & Verification & Status \\
\midrule\endhead
\bottomrule\endfoot
\multicolumn{5}{@{}l}{\textit{Frame engine and voice management}} \\
CT1 & SYNC at \SI{1}{kHz} from the \SI{8}{MHz} crystal on TIM1 CH1 (PA8). At each edge: FRAME, then enumeration
  packets, queued packets up to 95 bytes, refreshes, all in one DMA transfer that has started within
  \SI{40}{\micro s} of the edge. SYNC jitter stays unchanged by USB traffic. & P1, P10, bus.md 3 & B (\texttt{usbtest.py jitter}: 0.6--0.9\,ns) & \stmet \\
CT2 & At start-up: all slots held in reset through the 74HC595, then released slot by slot. ASSIGN and HELLO each
  slot; a slot without HELLO within \SI{300}{ms} is empty. CONFIG and the patches go to all. The chip ID is mapped
  to its slot. & F1, P9, backplane.md 3.5 & T \texttt{reset}; H (V3) & \stpart{} (one slot without NRST; the 595
  path is not written) \\
CT3 & Supervision: the controller polls STATUS. A card that disappears or resets is found and enumerated again;
  a card that hangs is reset through NRST, or over SWIO if that fails. & bus.md 6 & T \texttt{reset} & \stpart{}
  (periodic ASSIGN on the bench) \\
CT4 & Voice allocation within a part's slot mask: a free voice, else the oldest released, else the oldest held.
  PATCH\_SELECT goes in the NOTE\_ON's frame when a voice changes parts; PATCH\_DATA goes first if the patch is not
  in the card's slots. & F6, midi.md 2 & T (to write) & \stpart{} (allocator over 8 voices, one patch) \\
CT5 & Scheduling: every event is stamped with frame $a + L$, and notes go ahead of controllers. Queued controller
  values of the same voice and ID are coalesced. An event that leaves after its frame is counted. & bus.md 3.2, 9.5
  & T \texttt{load} & \stpart{} (no coalescing) \\
CT6 & Refreshes: VOICE\_STATE of every voice within \SI{50}{ms}; the active patches round robin through PATCH\_DATA;
  a refresh never overtakes a queued note. & bus.md 5.3 & T \texttt{repair} & \stpart{} (PATCH\_PRESET) \\
\multicolumn{5}{@{}l}{\textit{MIDI and parts}} \\
CT7 & A class-compliant USB-MIDI device (one cable each way) on the backplane's USB-C, which is USBD (PA11 / PA12).
  Flow control by NAK: no message is lost. & F2 & T \texttt{usbtest.py all}; H (USBD) & \stpart{} (shown on USBFS
  with a computer; USBD enumerated by the V203's own host) \\
CT8 & USB host for class-compliant MIDI keyboards on USBFS (PB6 / PB7): power switch HOST\_EN, fault input
  HOST\_FLT\_N. & F3 & T \texttt{hosttest.py}; B (development board: its USB-C is USBFS); H & \stpart{}
  (self-loop: enumeration, MIDI to the card, detach; no real keyboard yet) \\
CT9 & MIDI implementation of midi.md 3: notes and velocity; pitch bend; channel pressure; program change;
  CC 1, 5, 7, 11, 64 (sustain), 65, 71--78; NRPN over all patch IDs; CC 120, 121, 123; MIDI clock, start and stop.
  & F2, F6 & T (to write) & \stpart{} (notes, program change, bend, CC 1 / 7 / 11 / 71 / 74) \\
CT10 & Up to 8 parts, each with: channel; key range; slot mask; mode (poly, mono, legato, unison, kit); patch;
  volume; transpose; pan. Default setup: part 1 on channel 1, all voices, poly. & F6, F7 & T (to write) &
  \stpart{} (one poly part, or one part per channel) \\
CT11 & Kit mode: notes map to (patch, pitch), using up to 4 patches preloaded into the cards of the kit's mask. &
  midi.md 2.1 & T (to write) & \stopen \\
\multicolumn{5}{@{}l}{\textit{Configuration, storage and updates}} \\
CT12 & Patch library of 64 programs (144 bytes each) and the setups, in the V203's flash with wear-levelled pages;
  working copies in SRAM; an empty or corrupt store falls back to the built-in default. & F6, midi.md 1 & T (to
  write) & \stopen \\
CT13 & CDC command interface (text, one command per line): patches, setups, parts, card status, calibration,
  firmware update, a bridge to one card's console. SysEx (ID 0x7D) offers the same patch and setup dumps; the
  universal identity request is answered. & F4, midi.md 6 & T (to write) & \stpart{} (binary bench commands;
  identity reply) \\
CT14 & Card firmware update: an image is received over CDC and programmed into one card or all over that slot's SWIO
  line (8 lines, PA2--PA7, PB12, PB13), then verified; the card is enumerated again. & F4 & B (Ardulink bridge on
  PB13) & \stpart{} (bridge only; on-controller programming open) \\
CT15 & Calibration: CALIBRATE per card, CAL\_DATA read back and stored by chip ID; sent back after a card has moved
  slots. & F9, bus.md 6 & H & \stopen \\
CT16 & Controller firmware update in the field without a programmer: the factory bootloader through BOOT0 (JP2) over
  USB or USART1 (J6). & F4 & H (or the development board) & \stopen{} (a software jump was inconclusive) \\
\multicolumn{5}{@{}l}{\textit{Options, resources and robustness}} \\
CT17 & USB audio input (UAC1) of the mix through the ADC taps (PA0, PA1). & F10 & H & \stopen{} (option) \\
CT18 & CLK\_OUT: \SI{24.000}{MHz} on TIM3 CH4 (PB1) for the cards' CLK\_IN, off by default. & spec 6 & H &
  \stopen{} (option) \\
CT19 & The USB D+ pull-up is enabled only while VBUS\_DET (PB15) sees a host. & controller.md 1 & H & \stopen \\
CT20 & Resources: \SI{64}{KB} flash and \SI{20}{KB} SRAM, USB stacks included; at least \SI{2}{KB} of stack never
  used (\emph{proposed}). & C3 & I (link map), B & \stmet{} for the bench build (\SI{20.5}{KB}, \SI{8.0}{KB}) \\
CT21 & Robustness: a watchdog. After any controller reset the cards keep running muted until they are reset and
  enumerated again; nothing depends on the cards staying in reset while the controller boots. & backplane.md 3.5 &
  B; H & \stdes \\
CT22 & Effects slot: an effects card in slot~1 is recognised at enumeration (HELLO capability bit), excluded from
  voice allocation, and receives its parameters as timed messages; MIDI CC~91 sets the reverb level. &
  Section~\ref{sec:fxslot} & H (with an effects card) & \stopen{} (backplane prepared) \\
\end{longtable}
\endgroup

% ---------------------------------------------------------------------------------------------------------------------
\subsection{Tools and non-functional requirements}\label{sec:srstl}

\begingroup\small\setlength{\tabcolsep}{4pt}
\begin{longtable}{@{}>{\raggedright}p{0.9cm} >{\raggedright}p{7.5cm} >{\raggedright}p{1.5cm}
  >{\raggedright}p{2.7cm} >{\raggedright\arraybackslash}p{2.6cm}@{}}
\caption{Tools and non-functional requirements.}\label{tab:srstl}\\
\toprule
ID & Requirement & Source & Verification & Status \\
\midrule\endfirsthead
\toprule
ID & Requirement & Source & Verification & Status \\
\midrule\endhead
\bottomrule\endfoot
TL1 & A bench test suite runs unattended against one card and up to 7 virtual voices: \texttt{bustest.py all} (17
  checks), \texttt{stress}, \texttt{load}, \texttt{latency}, \texttt{counters}; \texttt{usbtest.py all}. & C5 & I
  & \stmet \\
TL2 & A host test compares the DAC renderer, compiled for x86, with a reference revision to the duty code; on the card,
  a console command prints a hash of the same render. & CF11 & T \texttt{render\_test.py} & \stmet \\
TL3 & Tables in the firmware (presets, CV laws, mode tables) are generated from the result files by
  \file{gen\_tables.py}. The generator checks the invariants, among them that a patch's oscillator levels sum to at
  most 1. & C5 & I & \stmet \\
TL4 & Environment-specific settings (lab host, ports, programmer serials) live only in the gitignored
  \file{tools/lab.conf}; the repository contains no host names, user names or local paths. & -- & I & \stmet \\
NF1 & Bare metal on ch32fun: no RTOS, no dynamic memory, all buffers static; interrupt priorities and nesting set
  explicitly. & C5 & I & \stmet \\
NF2 & Builds with \texttt{make} in WSL (riscv64-unknown-elf-gcc); the build variants of Table~\ref{tab:sadbuild}
  build without errors. & C5 & B & \stmet \\
NF3 & Every timing-critical path has a measured budget and a counter visible over the bus or the console: tick run
  time, render time and slack, missed halves, late events, stack low-water mark. & -- & T \texttt{counters} &
  \stmet \\
NF4 & Forward compatibility: HELLO carries the firmware version and capabilities; unknown message types and PARAM IDs
  are ignored and counted, never fatal. & bus.md 4 & T \texttt{ctrl} & \stmet \\
\end{longtable}
\endgroup

% ---------------------------------------------------------------------------------------------------------------------
\subsection{Software checks before the boards are ordered}\label{sec:srsorder}

What has to hold before the PCBs are ordered is not the software's completeness. It is whether the boards fix
anything the software cannot change later: which pin drives what, which timer and DMA channel serve it, and what the
driver circuits need. Table~\ref{tab:srsorder} lists these checks; the four marked \emph{new} were done on
20 October 2026, three of them on the bench V006.

\begingroup\small\setlength{\tabcolsep}{4pt}
\begin{longtable}{@{}>{\raggedright}p{5.4cm} >{\raggedright}p{4.6cm} >{\raggedright\arraybackslash}p{5.6cm}@{}}
\caption{Hardware commitments the software depends on.}\label{tab:srsorder}\\
\toprule
Check & Fixed by the boards & Result \\
\midrule\endfirsthead
\toprule
Check & Fixed by the boards & Result \\
\midrule\endhead
\bottomrule\endfoot
\multicolumn{3}{@{}l}{\textit{Done}} \\
Card pin map in the firmware (\file{board.h}) against the schematic netlist (\file{hw/voice.py}) & all 20 MCU pins
  & agree pin for pin (inspection, \emph{new}) \\
V006 reset pin default & no NRST series resistor or pull-down; 595 drives NRST push-pull & option byte USER = 0xF7:
  RST\_MODE = 10b, PD7 is NRST with a \SI{12}{ms} filter, the factory default (reference manual 18.5; read from the
  bench chip, \emph{new}) \\
TIM1 remap 0010b on the V006: CV PWM pins & CV\_RES, CV\_CUT, VCA trim, CV\_VCA on PC0, PC7, PC6, PD3 & all four
  pins toggle with a note held (GPIO input registers sampled, \emph{new}) \\
TIM2 CH1 on PC5 beside the running DAC (pan) & PWM\_PAN on PC5 & PC5 toggles at the carrier with CH1 enabled by hand,
  the DAC unaffected (\emph{new}); firmware still to write (CF14) \\
ADC channel of the loopback & ADC\_AUDIO on PC4 & PC4 is ADC A2 (datasheet pin table); capture to be shown on a card
  \\
Bus pins of both MCUs & SYNC PA8 / PD4, BCAST PB10 / PD6, RETURN PB11 / PD5, SWIO PB13 / PD1 & bench, since
  2026-10-19 (the backplane uses the same V203 pins; PB13 is SWIO\_8) \\
DAC on TIM2 CH2 (PC2), DMA ch1 & PWM\_DAC & bench (carrier on PC2) \\
USB device without an external pull-up (USBFS) & -- (the host port) & enumerates on the development board \\
\multicolumn{3}{@{}l}{\textit{Recommended before ordering (a resistor and a wire on the bench)}} \\
Return line open drain with the backplane's \SI{1}{k\ohm} pull-up: the default card build (no \texttt{BUS\_TX\_PP})
  & R on RETURN (\SI{1}{k\ohm}), card TX open drain & \stopen: fit \SI{1}{k\ohm} from PB11 to 3.3\,V, build without
  \texttt{BUS\_TX\_PP}, run \texttt{bustest.py all} \\
NRST driven push-pull, as by the 74HC595 & U6 drives each slot's NRST directly & \stopen: wire a V203 GPIO to the
  V006's NRST, reset and re-enumerate (\texttt{bustest.py reset} with NRST) \\
\multicolumn{3}{@{}l}{\textit{Can follow the order (footprints or jumpers cover them)}} \\
USBD (PA11 / PA12) enumeration and the internal D+ pull-up & the \SI{1.5}{k\ohm} footprint (not fitted) &
  enumerated on the internal pull-up against the V203's own host (\emph{new}); a computer as the host still to try \\
USBFS host pull-downs & two \SI{15}{k\ohm} footprints (not fitted) & internal, about 13--15\,k$\Omega$ (\emph{new}) \\
Factory bootloader through BOOT0 & JP2, J6 (USART1) & USART1 and SWD remain if the bootloader's USB port is the
  other one \\
74HC595 at \SI{9}{MHz} and \SI{3.3}{V}; SPI1 remap & U6 on SPI1 remap 1 & the clock can be lowered in firmware \\
CLK\_IN \SI{24}{MHz}, SPI flash, MIDI DIN & not fitted & options \\
\end{longtable}
\endgroup

\paragraph{Assessment.} Nothing found so far blocks the order. The two recommended checks take one resistor and one
wire on the bench and exercise the production configuration of the return line and the reset path. Both have so
far only run in bench variants. The firmware still to write (pan, PATCH\_DATA and the full parameter set, the
\SI{5}{V} calibration, the backplane controller) depends on these pins and circuits, but none of it changes them.

% ---------------------------------------------------------------------------------------------------------------------
\subsection{Decisions for review}\label{sec:srsdec}

\begin{enumerate}
  \item \textbf{Proposed margins.} CF12 (50 carrier periods of slack), CF21 and CT20 (RAM, stack, flash limits) are
    proposals and come from no system requirement.
  \item \textbf{Pan ownership.} CF14 puts the pan law (constant power, the CV\_VCA cap) in the card and the position
    in a per-voice PARAM. The alternative is to let the controller compute both and send raw values.
  \item \textbf{Parameter coverage for the first build.} Is the full patch set of midi.md 4 needed (CF15), or the
    subset of the built-in presets plus the most used IDs, with PATCH\_PRESET kept?
  \item \textbf{CDC format.} CT13 asks for text commands (as the board~1 console). The bench uses binary commands,
    which the test suite relies on. Keep both?
  \item \textbf{Storage.} CT12 uses the V203's internal flash (about \SI{12}{KB} of 64). The SPI flash footprint (not
    fitted) would hold more patches or samples later.
  \item \textbf{Controller updates.} If the factory bootloader does not run on the USB-C port (CT16), is USART1 (J6)
    plus SWD enough for field updates, or should the controller get an own USB update path (a DFU-like mode in its
    firmware)?
  \item \textbf{Board~1.} Keep board~1 (the rev 0.1 test bed, V003 pin functions plus its DCO) in the code base, or
    freeze it at the current revision?
\end{enumerate}
